% The split protocol is frozen (AGENTS.md section 2). It changes only through a
% decision record; this chapter describes it, it does not redefine it.
\chapter{Experimental Setup}
\label{ch:experiments}

\section{Splits}
\label{sec:splits}
% source: docs/decisions/0006-floorplancad-split-by-source-group.md
The public-data split keeps related drawings together by assigning complete source
groups to training, validation or held-out evaluation, with each source identified
by its numbering namespace and file-name prefix rather than by the prefix alone.
% source: docs/decisions/0006-floorplancad-split-by-source-group.md
The configured namespaces distinguish the separately numbered release folder from
the folders that share a numbering scheme.
% source: docs/decisions/0006-floorplancad-split-by-source-group.md
Sources containing a byte-identical file are merged into the same group, and exact
copies are removed before the split is constructed.
% source: docs/decisions/0007-floorplancad-fold-assignment.md
% source: docs/decisions/0006-floorplancad-split-by-source-group.md
Where copies share identical content, the copy with the smallest drawing identifier
is retained, and every resulting group belongs entirely to a single fold.

% source: docs/decisions/0006-floorplancad-split-by-source-group.md
The official release folders are not used as evaluation folds, because the recorded
split audit found source groups and identical content shared between the official
training and test folders.
% source: docs/decisions/0006-floorplancad-split-by-source-group.md
The earlier assumption that prefixes identify a common source across all release
folders was also rejected: the numbering namespace is needed to interpret them.
% source: docs/decisions/0006-floorplancad-split-by-source-group.md
These checks establish the grouping rule used here.

% source: docs/decisions/0007-floorplancad-fold-assignment.md
Groups are sorted by identifier and shuffled using the seed stored in the split
configuration.
% source: docs/decisions/0007-floorplancad-fold-assignment.md
Complete groups are assigned to the held-out fold until its target fraction of
deduplicated drawings is reached, after which the validation fold is filled in the
same way and the remaining groups form the training fold.
% source: docs/decisions/0007-floorplancad-fold-assignment.md
The fractions apply to drawings rather than groups because group sizes vary, and
keeping a group intact can cause a fold to exceed its target size.

% source: docs/decisions/0007-floorplancad-fold-assignment.md
The resulting assignment is stored in \texttt{split.csv} under its split run
identifier, and subsequent experiments load that assignment instead of recomputing
it.
% source: docs/decisions/0007-floorplancad-fold-assignment.md
This keeps the evaluation drawings fixed when the splitting code changes.
% source: docs/decisions/0007-floorplancad-fold-assignment.md
Classes are not stratified; their distribution within each fold is to be reported
alongside the split.

% source: docs/decisions/0006-floorplancad-split-by-source-group.md
Exact deduplication does not detect overlapping crops with different content, so
near-duplicate drawings across source groups remain an unresolved limitation.
% source: docs/decisions/0007-floorplancad-fold-assignment.md
Whether class imbalance between folds requires a revised assignment is also left
open in the decision record.
% source: docs/decisions/0007-floorplancad-fold-assignment.md
A revised split requires a new decision and a new run.

% source: runs/2026-10-26-floorplancad-split/metrics.json
The recorded fold sizes and target-class counts are reported in
\cref{tab:public-folds,tab:public-class-counts}.
% source: runs/2026-10-26-floorplancad-split/metrics.json
The inventory column counts all parsed entities, including text, whereas the
target-class columns report the mapped wall, door and window counts.
% source: docs/decisions/0009-element-metrics.md
Those inventory totals are descriptive quantities, rather than the denominator
of an entity-classification performance score.

\begin{table}[htbp]
  \centering
  % source: runs/2026-10-26-floorplancad-split/metrics.json
  \caption{Frozen public-data fold assignment. Parsed entity totals include text;
  only the geometric entity types specified in the metric protocol are scored.}
  \label{tab:public-folds}
  \input{tables/public_folds.tex}
\end{table}

\begin{table}[htbp]
  \centering
  % source: runs/2026-10-26-floorplancad-split/metrics.json
  \caption{Target-class entity counts in the frozen folds, using the adopted
  release-specific mapping.}
  \label{tab:public-class-counts}
  \input{tables/public_class_counts.tex}
\end{table}

\section{Metrics}
\label{sec:metrics}
% source: docs/decisions/0009-element-metrics.md
Evaluation is performed on geometric entities, with precision, recall and the
harmonic-mean \(F\)-score computed for each target class; walls are the primary
target, while doors and windows are reported separately.
% source: docs/decisions/0009-element-metrics.md
Text entities are excluded from scoring because they carry no semantic labels in
the downloaded release.
% source: docs/decisions/0009-element-metrics.md
The scored entity types are lines, arcs, circles and ellipses, with each reference
class derived from the configured mapping of semantic identifiers
to wall, door or window; identifiers outside those lists and missing labels are
assigned to \texttt{other}.
% source: docs/decisions/0009-element-metrics.md
Every evaluated model must produce a class prediction for each scored entity.

% source: docs/decisions/0009-element-metrics.md
For a given target class, a true positive is an entity whose prediction and
reference label both belong to that class.
% source: docs/decisions/0009-element-metrics.md
A false positive is an entity predicted as that class whose reference label differs,
and a false negative is an entity belonging to that class whose prediction differs.
% source: docs/decisions/0009-element-metrics.md
Precision divides true positives by all positive predictions, while recall divides
true positives by all reference positives, and the \(F\)-score is their harmonic mean.

% source: docs/decisions/0009-element-metrics.md
Counts are pooled over the held-out fold before the class scores are computed,
giving each entity equal weight rather than giving each drawing equal weight.
% source: docs/decisions/0009-element-metrics.md
A score with an empty denominator is assigned zero under the adopted convention,
and per-drawing scores may supplement the pooled scores but do not replace them.
% source: docs/decisions/0009-element-metrics.md
Raster predictions are evaluated only after their pixel outputs have been mapped
back to the scored vector entities.
% source: docs/decisions/0009-element-metrics.md
The representation comparison therefore requires the mapping described in
\cref{sec:raster-to-vector} before raster scores can be included, so that both
branches are compared using the same unit of evaluation.

\section{Implementation and Environment}
\label{sec:implementation}
% source: docs/decisions/0008-floorplancad-only-training.md
The planned external evaluation constrains the input representation: features and
models may use only information that the conversion from industrial drawings can
reproduce.
% source: docs/decisions/0008-floorplancad-only-training.md
The recorded constraints include compatible entity kinds, crop size and layer groups.
% source: docs/decisions/0008-floorplancad-only-training.md
The conversion procedure remains to be specified.
% source: docs/decisions/0005-floorplancad-class-mapping.md
% source: docs/decisions/0007-floorplancad-fold-assignment.md
The class mapping, split seed and target fold fractions are configuration settings, while
the actual drawing assignments are loaded from the recorded split run.
% source: docs/decisions/0007-floorplancad-fold-assignment.md
Experiments cite that run identifier to identify their evaluation partition.

\todo[inline]{Add the selected model architectures, geometry and layer-name inputs,
preprocessing, training settings, model-selection procedure and stopping criteria
from implementation decisions and training runs. Define the geometry-only
comparison explicitly before claiming that it answers the layer-name question.}
\todo[inline]{Document the pinned software environment, hardware and code snapshot from
the corresponding run records. Training execution and reproducibility have not yet
been established by the decisions cited in this chapter.}
